%% ============================================================
%%  Project Report — Machine Learning, Master 203, Dauphine
%%  Iterated Prisoner's Dilemma & Q-Learning
%% ============================================================

\documentclass[11pt, a4paper]{article}

%% ---- Encoding & language ----
\usepackage[utf8]{inputenc}
\usepackage[T1]{fontenc}
\usepackage[english]{babel}

%% ---- Page layout ----
\usepackage{geometry}
\geometry{top=2.5cm, bottom=2.5cm, left=2.5cm, right=2.5cm}
\usepackage{setspace}
\onehalfspacing

%% ---- Typography ----
\usepackage{lmodern}
\usepackage{microtype}

%% ---- Mathematics ----
\usepackage{amsmath, amssymb, amsthm}
\usepackage{mathtools}

%% ---- Figures & tables ----
\usepackage{graphicx}
\usepackage{booktabs}
\usepackage{array}
\usepackage{caption}
\usepackage{subcaption}
\usepackage{float}

%% ---- Colours & code ----
\usepackage{xcolor}
\usepackage{listings}
\lstset{
  basicstyle=\footnotesize\ttfamily,
  keywordstyle=\color{blue},
  commentstyle=\color{gray},
  breaklines=true,
  frame=single,
  numbers=left,
  numberstyle=\tiny\color{gray}
}

%% ---- Hyperlinks ----
\usepackage{hyperref}
\hypersetup{
  colorlinks=true,
  linkcolor=blue!70!black,
  citecolor=green!50!black,
  urlcolor=blue!80!black
}

%% ---- Bibliography ----
\usepackage{natbib}
\bibliographystyle{plainnat}

%% ---- Headers & footers ----
\usepackage{fancyhdr}
\pagestyle{fancy}
\fancyhf{}
\rhead{Master 203 — Machine Learning}
\lhead{Prisoner's Dilemma \& Q-Learning}
\cfoot{\thepage}

%% ---- Mathematical environments ----
\newtheorem{definition}{Definition}[section]
\newtheorem{proposition}{Proposition}[section]
\newtheorem{remark}{Remark}[section]

%% ============================================================
%%  METADATA
%% ============================================================

\title{
  \vspace{-1cm}
  {\Large Project Report — Machine Learning}\\[0.4em]
  {\LARGE \textbf{Reinforcement Learning in the Iterated Prisoner's Dilemma:}}\\[0.3em]
  {\Large Convergence and Comparative Analysis of Strategies}
}

\author{
  Gabriel \textsc{Tseng} \and
  Axel \textsc{Rubio} \and
  Benjamin \textsc{Krief} \and
  Krea \textsc{Fares}\\[0.5em]
  \small Master 203 — Université Paris-Dauphine\\
  \small Machine Learning Course — \today
}

\date{}
\begin{document}

\maketitle
\thispagestyle{empty}

%% ---- Abstract ----
\begin{abstract}

Game theory provides a rigorous framework for analysing strategic interactions in which individual incentives diverge from collective welfare, 
a tension that arises naturally in multi-agent systems like financial markets. 
The Iterated Prisoner's Dilemma (IPD) is a canonical instance of this tension, and applying reinforcement learning to it allows us to study how adaptive agents 
learn to navigate without any prior knowledge of the game structure.\\

This report will go through the application of Q-learning to the IPD. 
There are two objectives: first, to understand how a rational learning agent behaves when facing classical deterministic strategies; 
second, to assess the relative importance of Q-learning hyperparameters : learning rate $\alpha$, discount factor $\gamma$, $\varepsilon$-greedy exploration policy.\\

Our experiments show that a Q-learning agent converges to systematic defection when facing any player that do not punish defection, consistent with the dominant Nash equilibrium of the one-shot game. 
We compare its performance against a set of baseline strategies, ranging from naive heuristics (\textit{Always Cooperate}, \textit{Always Defect}) to reactive strategies (\textit{Tit for Tat} and its variants), and discuss the conditions under which cooperation can or cannot emerge in a mixed population.
\end{abstract}

\newpage
\tableofcontents
\newpage

\section{Introduction: The Prisoner's Dilemma}
\label{sec:intro}
%% ============================================================

\subsection{Game description}

The Prisoner's Dilemma is a symmetric two-player game popularised in the context of game theory by Robert Axelrod in 1984. At each round, two players simultaneously and independently choose between two actions: \textbf{cooperate} (C) or \textbf{betray} (B). The payoff matrix is defined as follows:

\begin{table}[H]
  \centering
  \caption{Payoff matrix of the Prisoner's Dilemma (row = player 1, column = player 2).}
  \label{tab:payoff}
  \begin{tabular}{c|cc}
    \toprule
                   & \textbf{Cooperate (C)} & \textbf{Betray (B)} \\
    \midrule
    \textbf{Cooperate (C)} & $(3,\, 3)$ & $(0,\, 4)$ \\
    \textbf{Betray (B)}    & $(4,\, 0)$ & $(1,\, 1)$ \\
    \bottomrule
  \end{tabular}
\end{table}

Those payoffs are defined such that they satisfy the following inequalities (as long as they do, the game is a Prisoner's Dilemma):

\begin{equation}
    T > R > P > S
\end{equation}

where

\begin{itemize}
    \item $T$ is the temptation to defect (4)
    \item $R$ is the reward for mutual cooperation (3)
    \item $P$ is the punishment for mutual defection (1)
    \item $S$ is the sucker's payoff (0)
\end{itemize}

The structure of this game gives rise to a fundamental tension between individual rationality and collective optimum: \textbf{betraying is a strictly dominant strategy} as it yields a higher payoff regardless of the other player's action, yet if both players choose it, the outcome $(B, B) = (1,1)$ is Pareto-dominated by $(C, C) = (3,3)$.

\end{document}